\documentclass[11pt,a4paper]{letter}

\usepackage[T1]{fontenc}
\usepackage[utf8]{inputenc}
\usepackage{lmodern}
\usepackage[a4paper,margin=2.5cm]{geometry}
\usepackage[hidelinks]{hyperref}
\setlength{\parskip}{0.8em}
\setlength{\parindent}{0pt}

\signature{Florian Maillard\\Independent researcher\\Strasbourg, France\\\href{mailto:florian.maillard@mailoo.org}{florian.maillard@mailoo.org}\\ORCID: 0000-0002-3334-725X}
\address{Florian Maillard\\Independent researcher\\Strasbourg, France}

\begin{document}

\begin{letter}{Editor\\Journal of Solid State Electrochemistry}

\opening{Dear Editor,}

Please find enclosed my manuscript entitled
\emph{Global resistance scaling and voltage-dependent local dynamics in Li-ion battery discharge}
for consideration for publication in the \emph{Journal of Solid State Electrochemistry}.

The manuscript examines ten archived lithium-ion battery discharges obtained
under two constant resistive loads, 3.33~$\Omega$ and 5~$\Omega$, with five
repeated records per condition. A minimal quasi-static one-state description
predicts the pointwise scaling
$\dot V(V,R_L)=G(V)/R_L$. The experimental analysis shows that this ideal
$1/R_L$ dependence is recovered very closely at the global level
($\alpha=1.498019$ compared with $5/3.33=1.501502$), while the corresponding
local voltage-rate ratio varies strongly with terminal voltage.

The main contribution is therefore a model-validity result rather than the
resistance ratio itself: a simple resistive scaling accurately captures the
global discharge-rate relation, but it is insufficient as a local description.
The local departure is robust to the tested smoothing parameters and bootstrap
resampling. Because the two resistance conditions were acquired sequentially,
the manuscript explicitly limits the causal interpretation of the local
difference and does not attribute it uniquely to resistance.

The measurements come from an experimental campaign that has previously been
used to study electrochemical voltage noise and higher-order fluctuation
statistics. The present manuscript addresses a distinct question: the slow
discharge trajectory and its voltage-rate scaling. It does not reuse the
spectral, harmonic, or fast-noise analyses of those earlier studies.

I believe the work is relevant to the journal because it connects a simple
electrochemical description of lithium-ion discharge with a direct experimental
test of its range of validity, while keeping the interpretation consistent with
the limitations of the archived experiment.

The manuscript includes the data and code availability statement, funding and
competing-interest declarations, and a statement describing the use of
AI-assisted tools. All analyses and interpretations were verified by the author.

Thank you for considering this manuscript.

\closing{Sincerely,}

\end{letter}
\end{document}
